\section{Attention and Weight Diagnostics}
\label{app:attn}

This appendix reports the mechanism measurements of Section \ref{sec:mechanism}: attention entropy and mass distribution in PatchTST, DLinear weight-decay profiles, and the repair ablation (scripts and per-run data are released in the accompanying repository).

\subsection{PatchTST attention entropy and mass distribution}

We load the trained PatchTST checkpoints ($d_{\mathrm{model}}{=}512$, $e_{\mathrm{layers}}{=}2$, $d_{\mathrm{ff}}{=}2048$) and, in eval mode (dropout off) over the first 8 test batches, record per layer: the per-query mean attention entropy $H$ (also normalized by $\ln S$), the attention mass on the first (oldest) patch (``sink''), and the total mass on the most recent 10\% of patches. Patch count is $S = (L-16)/8 + 2$. Reference points: uniform attention gives $H/\ln S = 1$, recency mass $0.10$, and sink mass $1/S$.

\begin{table}[t]
\caption{Exchange rate, attention statistics per layer (L0/L1 = first/second layer), mean over seeds 2021/2022. Entropy approaches the uniform ceiling as $L$ grows and the sink dissolves from $0.138$ (twice the uniform level $1/S$ at $L{=}96$) to the uniform level $1/S \approx 0.003$, but the recency mass is never diluted below its uniform share.}
\label{tab:attn_exchange}
\centering\small
\begin{tabular}{cc cccc cc cc}
\toprule
$L$ & $S$ & $H$(L0) & $H/\ln S$(L0) & $H$(L1) & $H/\ln S$(L1) & sink L0 & sink L1 & near-10\% L0 & near-10\% L1 \\
\midrule
96 & 12 & 1.92 & 0.77 & 2.09 & 0.84 & 0.138 & 0.127 & 0.127 & 0.103 \\
720 & 90 & 4.18 & 0.93 & 4.07 & 0.90 & 0.021 & 0.015 & 0.124 & 0.087 \\
2880 & 360 & 5.19 & 0.88 & 5.38 & 0.91 & 0.007 & 0.004 & 0.163 & 0.145 \\
\bottomrule
\end{tabular}
\end{table}

\begin{table}[t]
\caption{Electricity, attention statistics, measured on fully trained uncapped vanilla control runs (Appendix repair ablation, \texttt{none} variant), since the capped $L{\ge}1440$ checkpoints of the main matrix were lost to OOM. Dilution is \emph{stronger} than on exchange rate ($H/\ln S$ at L1 reaches 0.95), yet electricity is where long context helps most.}
\label{tab:attn_electricity}
\centering\small
\begin{tabular}{cc cc cc cc}
\toprule
$L$ & $S$ & $H/\ln S$(L0) & $H/\ln S$(L1) & sink L0 & sink L1 & near-10\% L0 & near-10\% L1 \\
\midrule
96 & 12 & 0.59 & 0.66 & 0.086 & 0.080 & 0.063 & 0.063 \\
1440 & 180 & 0.66 & 0.91 & 0.008 & 0.005 & 0.185 & 0.093 \\
2880 & 360 & 0.75 & 0.95 & 0.003 & 0.003 & 0.192 & 0.129 \\
\bottomrule
\end{tabular}
\end{table}

On the capped main-matrix checkpoints, electricity at $L{=}96$ shows $H/\ln S = 0.60/0.68$ and near-10\% mass $0.063/0.068$ (below the uniform $0.083$); a partially trained $L{=}720$ checkpoint shows $0.77/0.88$ and $0.190/0.114$. Two findings follow. First, softmax dilution is real and universal: normalized entropy rises toward the uniform ceiling on both datasets, and no attention sink survives at long lookbacks. Second, dilution is not harmful by itself: it is stronger on electricity, which benefits from long context, and the recency mass is not diluted on either dataset (exchange $0.13 \to 0.16$, electricity L0 $0.06 \to 0.19$). Whether flattening hurts depends on whether the diluted far content is noise (exchange) or periodic structure (electricity) --- consistent with the DLinear evidence next.

\subsection{DLinear weight profiles}

DLinear learns a linear map per decomposition branch; averaging $|w|$ over output positions gives a mass profile over lag positions (index $0$ = oldest). Table \ref{tab:dlinear} reports, at $L{=}2880$: the mass on the most recent 96 lags, the mass beyond lag 720, and the daily-harmonic contrast (peak power at the daily period relative to its spectral neighborhood).

\begin{table}[t]
\caption{DLinear weight-mass distribution at $L{=}2880$. On exchange rate the seasonal branch places only $6.0\%$ of its mass on the recent 96 lags and $72.6\%$ beyond lag 720 with no periodic structure (contrast $\approx 1$): the linear model does not learn implicit gating, and the diffuse far-lag mass coincides with the $L{=}2880$ degradation (MSE $0.636$). On electricity the far mass is equally large but carries daily/weekly spikes (contrast $1.25$; at $L{=}96$ the seasonal profile shows sharp $d{\approx}24/168$ peaks with contrast $1.67$), and long context helps (MSE $0.129$).}
\label{tab:dlinear}
\centering\small
\begin{tabular}{llccc}
\toprule
Dataset & Branch & Mass(last 96 lags) & Mass(lag $> 720$) & Daily contrast \\
\midrule
Exchange & Linear\_Seasonal & 0.060 & 0.726 & 1.03 \\
Exchange & Linear\_Trend & 0.238 & 0.556 & 1.05 \\
Electricity & Linear\_Seasonal & 0.187 & 0.487 & 1.25 \\
Electricity & Linear\_Trend & 0.181 & 0.529 & 0.94 \\
\bottomrule
\end{tabular}
\end{table}

\subsection{Repair ablation (attention-side fixes)}

We test two attention-internal repairs on PatchTST (\texttt{PatchTSTGated}, an SDPA re-implementation whose \texttt{none} variant is numerically equal to PatchTST, smoke-tested at $\max|\Delta| = 9.5{\times}10^{-7}$): \texttt{denbias}, a learnable off-by-one null logit that lets softmax attend to ``nothing'', and \texttt{recmask}, a hard recency mask (identity at $L{=}96$, where the patch count $12$ is below the mask window). Protocol: exchange rate and electricity $\times$ $L \in \{96, 1440, 2880\}$ $\times$ variant $\times$ seeds $\{2021, 2022\}$, uncapped windows, main-matrix hyperparameters ($d_{\mathrm{model}}{=}512$, $d_{\mathrm{ff}}{=}2048$, $e_{\mathrm{layers}}{=}2$, 8 heads, LR $10^{-4}$; batch size 32 on exchange, $16/8/4$ on electricity for $L = 96/1440/2880$).

\begin{table}[t]
\caption{Repair ablation, test MSE (mean $\pm$ std over seeds 2021/2022). The hard recency mask recovers $34\%$ of the exchange-rate degradation at $L{=}2880$ ($0.399 \to 0.263$, closing $\sim$45\% of the gap to short lookbacks) without hurting electricity; the learnable softmax exit is never used (its null logit stays at $\approx 0$ and the null mass stays $\le 0.4\%$ on exchange), so \texttt{denbias} is pointwise identical to the control. Sanity checks: \texttt{none} at $L{=}2880$ reproduces the capped main-matrix run ($0.399$ vs.\ $0.405$; that run's window count equals the uncapped one), and \texttt{recmask} at $L{=}96$ is bitwise identical to \texttt{none}.}
\label{tab:repair}
\centering\small
\begin{tabular}{lcccc}
\toprule
Dataset / $L$ & none (control) & denbias & recmask & capped PatchTST \\
\midrule
Exchange / 96 & $0.0950\pm0.0009$ & $0.0952\pm0.0010$ & $0.0950\pm0.0009$ & $0.0871\pm0.0020$ \\
Exchange / 1440 & $0.1371\pm0.0003$ & $0.1387\pm0.0030$ & $0.1282\pm0.0132$ & $0.1484\pm0.0239$ \\
Exchange / 2880 & $0.3989\pm0.0168$ & $0.3989\pm0.0181$ & $\mathbf{0.2625\pm0.0424}$ & $0.4053\pm0.0169$ \\
\midrule
Electricity / 96 & $0.1803\pm0.0002$ & $0.1803\pm0.0002$ & $0.1803\pm0.0002$ & $0.1816\pm0.0001$ \\
Electricity / 1440 & $0.1291\pm0.0003$ & $0.1285\pm0.0003$ & $0.1291\pm0.0003$ & -- (OOM) \\
Electricity / 2880 & $0.1301\pm0.0009$ & $0.1384\pm0.0089$ & $\mathbf{0.1269\pm0.0000}$ & -- (OOM) \\
\bottomrule
\end{tabular}
\end{table}

About half of the long-context harm on exchange rate is directly attributable to attention being forced to read the entire far past (a hard mask recovers it); the remainder sits in the value/head pathway (on electricity at $L{=}2880$ the vanilla validation curve oscillates violently while \texttt{recmask}/\texttt{denbias} converge smoothly, indicating an optimization component). A learnable softmax exit is not an effective mechanism --- gradients never use it. (The \texttt{denbias} outlier on electricity at $L{=}2880$ is single-seed noise: seed 2022 early-stopped after epoch 1 at $0.1473$ while seed 2021 reached $0.1295$.)
